\newpage
% ---------------------------------------------------------------------
\subsubsection{Vincoli di Progettazione}
\label{subsec:vincoli-progettazione}
% ---------------------------------------------------------------------

Questa sezione descrive i vincoli di progettazione e le scelte architetturali e tecnologiche imposte dal committente o definite dal team per garantire il rispetto dei requisiti non funzionali, in particolare quelli di manutenibilità (REQ-NFR-05) e performance (REQ-NFR-01).

% ---------------------------------------------------------------------
\paragraph{Architettura Distribuita}
\label{subsubsec:architettura-distribuita}
% ---------------------------------------------------------------------

Il sistema \productname{} adotta un'architettura distribuita client-server con separazione netta tra presentazione e logica di business:

\begin{itemize}[leftmargin=1.1em]
\item \textbf{Front-end}: applicazione web Single Page Application (SPA) responsabile dell'interfaccia utente e dell'interazione
\item \textbf{Back-end}: server applicativo che gestisce logica di business, autenticazione, validazione dati e persistenza
\item \textbf{Database}: PostgreSQL per persistenza dati relazionali
\end{itemize}

La comunicazione tra front-end e back-end avviene esclusivamente tramite \textbf{API REST} su protocollo HTTP/HTTPS. Il front-end non ha accesso diretto al database; tutte le operazioni sui dati sono gestite dal back-end.

% ---------------------------------------------------------------------
\paragraph{Stack Tecnologico Back-end}
\label{subsubsec:stack-backend}
% ---------------------------------------------------------------------

\subparagraph{Tecnologie Core}

Il back-end utilizza il seguente stack tecnologico:

\begin{itemize}[leftmargin=1.1em]
\item \textbf{Linguaggio}: Java
\item \textbf{Framework}: Spring Boot
\item \textbf{Build Tool}: Apache Maven per gestione dipendenze e build automation
\item \textbf{Database Access}: Spring Data JPA
\item \textbf{Migrations}: Flyway per versioning dello schema database
\end{itemize}

\subparagraph{Moduli Spring Boot}

Il progetto integra i seguenti moduli dell'ecosistema Spring:

\begin{itemize}[leftmargin=1.1em]
\item \textbf{Spring Web}: per sviluppo API REST con annotazioni \texttt{@RestController} e serializzazione JSON automatica
\item \textbf{Spring Security}: per autenticazione JWT, password hashing con BCrypt, autorizzazione basata su ruoli
\end{itemize}

% ---------------------------------------------------------------------
\paragraph{Stack Tecnologico Front-end}
\label{subsubsec:stack-frontend}
% ---------------------------------------------------------------------

\subparagraph{Tecnologie Core}

Il front-end utilizza tecnologie moderne per lo sviluppo di applicazioni web:

\begin{itemize}[leftmargin=1.1em]
\item \textbf{Linguaggio}: TypeScript per type safety e migliore developer experience
\item \textbf{Framework UI}: React con architettura component-based
\item \textbf{Build Tool}: Vite per sviluppo rapido e build ottimizzate
\end{itemize}

\subparagraph{Styling e Animazioni}

L'interfaccia utente combina:

\begin{itemize}[leftmargin=1.1em]
\item \textbf{Tailwind CSS 3.x}: framework utility-first per styling rapido e consistente
\item \textbf{Framer Motion}: libreria per animazioni e transizioni fluide
\end{itemize}

Questa combinazione permette di creare un design system custom che emula esteticamente Material Design 3 pur mantenendo piena flessibilità di personalizzazione.

\subparagraph{HTTP Client e Gestione Stato}

\begin{itemize}[leftmargin=1.1em]
\item \textbf{Axios}: per chiamate API REST con interceptors per gestione automatica JWT
\item \textbf{Context API / Zustand}: per gestione stato globale dell'applicazione
\end{itemize}

% ---------------------------------------------------------------------
\paragraph{Database}
\label{subsubsec:database}
% ---------------------------------------------------------------------

\subparagraph{Sistema di Gestione Database}

Il sistema utilizza \textbf{PostgreSQL} come database relazionale per:
\begin{itemize}[leftmargin=1.1em]
\item Conformità ACID completa
\item Supporto transazioni e integrità referenziale
\end{itemize}

\subparagraph{Accesso ai Dati}

L'accesso al database avviene tramite:
\begin{itemize}[leftmargin=1.1em]
\item \textbf{Spring Data JPA}: per astrazione dell'accesso dati con repository pattern
\item \textbf{Hibernate}: come provider JPA per mapping object-relational
\item \textbf{Flyway}: per gestione evolution schema database tramite migration SQL versionate
\end{itemize}

% ---------------------------------------------------------------------
\paragraph{Deployment e Containerizzazione}
\label{subsubsec:deployment}
% ---------------------------------------------------------------------

\subparagraph{Docker}

Tutti i componenti del sistema sono containerizzati tramite \textbf{Docker} per garantire:
\begin{itemize}[leftmargin=1.1em]
\item Isolamento delle dipendenze
\item Portabilità tra ambienti (sviluppo, testing, produzione)
\item Riproducibilità delle configurazioni
\end{itemize}

Ogni componente ha un \textbf{Dockerfile} dedicato:
\begin{itemize}[leftmargin=1.1em]
\item Back-end: build multi-stage con Maven e runtime Java 17
\item Front-end: build con Node.js e serving con Nginx
\item Database: immagine ufficiale PostgreSQL
\end{itemize}

\subparagraph{Docker Compose}

L'orchestrazione dei container in ambiente di sviluppo è gestita tramite \textbf{Docker Compose}, che coordina:
\begin{itemize}[leftmargin=1.1em]
\item Container PostgreSQL con volume persistente
\item Container back-end con dipendenza dal database
\item Container front-end con dipendenza dal back-end
\item Network condivisa per comunicazione inter-container
\end{itemize}

% ---------------------------------------------------------------------
\paragraph{Sicurezza}
\label{subsubsec:sicurezza-sistema}
% ---------------------------------------------------------------------

\subparagraph{Autenticazione e Password}

\begin{itemize}[leftmargin=1.1em]
\item \textbf{Password hashing}: BCrypt tramite Spring Security
\item \textbf{Autenticazione}: JSON Web Tokens (JWT) con algoritmo HMAC SHA-26
\item \textbf{Protezione brute force}: blocco account per 15 minuti dopo 5 tentativi falliti
\end{itemize}

\subparagraph{Validazione}

\begin{itemize}[leftmargin=1.1em]
\item Validazione lato server obbligatoria per tutti gli input
\item Sanitizzazione dati per prevenzione SQL injection
\item Autorizzazione basata su ruoli verificata ad ogni richiesta API
\end{itemize}

% ---------------------------------------------------------------------
\paragraph{Ambiente di Sviluppo}
\label{subsubsec:ambiente-sviluppo}
% ---------------------------------------------------------------------

\subparagraph{Strumenti Consigliati}

\begin{itemize}[leftmargin=1.1em]
\item \textbf{IDE Backend}: IntelliJ IDEA o Eclipse con Spring Tools
\item \textbf{IDE Frontend}: Visual Studio Code con estensioni TypeScript, ESLint, Tailwind CSS IntelliSense
\item \textbf{Database Client}: pgAdmin o DBeaver per gestione PostgreSQL
\item \textbf{Version Control}: Git con hosting su GitHub o GitLab
\end{itemize}